\documentclass[11pt]{article}
\newcommand{\DocShort}{Revisi\'on por pares}
\usepackage{fontspec}
\usepackage[spanish,es-noshorthands]{babel}
\decimalpoint
\usepackage{geometry}
\geometry{letterpaper, top=2.5cm, bottom=2.5cm, left=2.8cm, right=2.8cm}
\usepackage[expansion=false]{microtype}
\usepackage{parskip}
\setlength{\parskip}{6pt}
\usepackage{amsmath,amssymb}
\usepackage{booktabs}
\usepackage{array}
\usepackage{longtable}
\usepackage{caption}
\usepackage[dvipsnames,table]{xcolor}
\usepackage{enumitem}
\setlist[itemize]{topsep=2pt, itemsep=2pt, parsep=0pt}
\setlist[enumerate]{topsep=2pt, itemsep=2pt, parsep=0pt}
\usepackage{fancyhdr}
\usepackage[hidelinks,pdfencoding=auto]{hyperref}
\usepackage{titlesec}
\usepackage{tcolorbox}
\tcbuselibrary{breakable}
\usepackage{pdfpages}

\definecolor{darkblue}{HTML}{1B3A5C}
\definecolor{medblue}{HTML}{2E6DA4}
\definecolor{lightblue}{HTML}{D6E8F7}
\definecolor{teal}{HTML}{1A7A6E}
\definecolor{lightteal}{HTML}{D0EDEA}
\definecolor{lightgrey}{HTML}{F5F5F5}
\definecolor{midgrey}{HTML}{6F6F6F}
\definecolor{warnyellow}{HTML}{FFF3CD}

\titleformat{\section}
{\large\bfseries\color{darkblue}}
{\thesection}{0.75em}{}[{\color{medblue}\titlerule[0.8pt]}]
\titleformat{\subsection}
{\normalsize\bfseries\color{darkblue}}
{\thesubsection}{0.75em}{}
\titleformat{\subsubsection}
{\normalsize\bfseries\color{medblue}}
{\thesubsubsection}{0.75em}{}

\pagestyle{fancy}
\fancyhf{}
\fancyhead[L]{\small\color{midgrey}ICV 442 Acueductos y Alcantarillados}
\fancyhead[R]{\small\color{midgrey}\DocShort}
\fancyfoot[C]{\small\color{midgrey}\thepage}
\renewcommand{\headrulewidth}{0.4pt}
\renewcommand{\footrulewidth}{0pt}

\rowcolors{2}{lightgrey}{white}
\renewcommand{\arraystretch}{1.18}

\newtcolorbox{infobox}{
  breakable,
  colback=lightblue,
  colframe=medblue,
  boxrule=0.8pt,
  arc=3pt,
  left=8pt,
  right=8pt,
  top=7pt,
  bottom=7pt
}

\newtcolorbox{warnbox}{
  breakable,
  colback=warnyellow,
  colframe=orange!70!black,
  boxrule=0.8pt,
  arc=3pt,
  left=8pt,
  right=8pt,
  top=7pt,
  bottom=7pt
}

\newcommand{\doccover}[3]{
  \begin{center}
  {\small\color{midgrey}Pontificia Universidad Cat\'olica Madre y Maestra\\Escuela de Ingenier\'ia Civil y Ambiental}

  \vspace{1.0cm}
  {\color{medblue}\rule{\linewidth}{2pt}}
  \vspace{0.3cm}

  {\Large\bfseries\color{darkblue}ICV-442-T Acueductos y Alcantarillado}\\[0.35em]
  {\LARGE\bfseries\color{darkblue}#1}\\[0.25em]
  {\large\itshape #2}

  \vspace{0.3cm}
  {\color{medblue}\rule{\linewidth}{2pt}}
  \vspace{0.9cm}
  \end{center}

  \begin{center}
  \begin{tabular}{ll}
  \toprule
  \textbf{Asignatura} & ICV-442-T Acueductos y Alcantarillado \\
  \textbf{Profesor} & Ricardo Rafael Hern\'andez Moreira, PhD \\
  \textbf{Periodo acad\'emico} & Mayo--agosto de 2026 \\
  \textbf{Proyecto} & Dise\~no de sistemas de saneamiento para Las Charcas, Azua \\
  \textbf{Versi\'on} & #3 \\
  \bottomrule
  \end{tabular}
  \end{center}

  \vspace{0.4cm}
}


\begin{document}
\doccover{Formulario de revisi\'on por pares}{Retroalimentaci\'on guiada entre equipos}{23 de mayo de 2026}

\begin{infobox}
Nombre sugerido del archivo:

\url{ICV442_T[EquipoRevisor]_RevisionPares_Eq[EquipoRevisado]_v1.0.pdf}

Si el espacio provisto no es suficiente, se pueden anexar p\'aginas adicionales con observaciones t\'ecnicas complementarias.
\end{infobox}

\section{Prop\'osito y uso}
\begin{itemize}
\item Este formulario se utiliza en la semana de revisi\'on por pares para documentar observaciones t\'ecnicas sobre el borrador de otro equipo.
\item La revisi\'on por pares es un requisito habilitante para la entrega final del proyecto.
\item La retroalimentaci\'on debe centrarse en consistencia t\'ecnica, claridad documental y verificabilidad, no en comentarios gen\'ericos.
\end{itemize}

\section{Identificaci\'on}
\begin{tabular}{p{0.32\linewidth} p{0.58\linewidth}}
\toprule
Equipo revisor & \\
Equipo revisado & \\
Fecha de revisi\'on & \\
Integrantes que revisan & \\
Documentos recibidos & \\
\bottomrule
\end{tabular}

\section{Verificaci\'on r\'apida del paquete revisado}
\begin{tabular}{p{0.54\linewidth} p{0.12\linewidth} p{0.20\linewidth}}
\toprule
\textbf{Aspecto} & \textbf{Estado} & \textbf{Observaci\'on breve} \\
\midrule
Se identifica claramente el alcance de los tres sistemas & S\'i / Parcial / No & \\
La memoria muestra criterios y supuestos legibles & S\'i / Parcial / No & \\
Los planos y esquemas son interpretables & S\'i / Parcial / No & \\
La nomenclatura de archivos es consistente & S\'i / Parcial / No & \\
Se distinguen modelos, anexos y soportes & S\'i / Parcial / No & \\
\bottomrule
\end{tabular}

\section{Revisi\'on t\'ecnica guiada}
\subsection*{1. Informaci\'on b\'asica y proyecci\'on}
\begin{tabular}{p{0.54\linewidth} p{0.12\linewidth} p{0.20\linewidth}}
\toprule
\textbf{Criterio} & \textbf{Estado} & \textbf{Observaci\'on} \\
\midrule
La fuente poblacional y el horizonte de dise\~no se entienden & S\'i / Parcial / No & \\
El m\'etodo de proyecci\'on parece justificable y trazable & S\'i / Parcial / No & \\
Las brechas de informaci\'on est\'an identificadas & S\'i / Parcial / No & \\
\bottomrule
\end{tabular}

\subsection*{2. Agua potable}
\begin{tabular}{p{0.54\linewidth} p{0.12\linewidth} p{0.20\linewidth}}
\toprule
\textbf{Criterio} & \textbf{Estado} & \textbf{Observaci\'on} \\
\midrule
Los caudales y par\'ametros principales se entienden & S\'i / Parcial / No & \\
La modelaci\'on EPANET parece consistente con el trazado & S\'i / Parcial / No & \\
Se reportan presiones y velocidades de forma verificable & S\'i / Parcial / No & \\
\bottomrule
\end{tabular}

\subsection*{3. Alcantarillado sanitario}
\begin{tabular}{p{0.54\linewidth} p{0.12\linewidth} p{0.20\linewidth}}
\toprule
\textbf{Criterio} & \textbf{Estado} & \textbf{Observaci\'on} \\
\midrule
Los caudales y criterios de dise\~no parecen claros & S\'i / Parcial / No & \\
Las pendientes, di\'ametros o profundidades parecen defendibles & S\'i / Parcial / No & \\
La modelaci\'on SWMM sanitaria se entiende & S\'i / Parcial / No & \\
\bottomrule
\end{tabular}

\subsection*{4. Alcantarillado pluvial}
\begin{tabular}{p{0.54\linewidth} p{0.12\linewidth} p{0.20\linewidth}}
\toprule
\textbf{Criterio} & \textbf{Estado} & \textbf{Observaci\'on} \\
\midrule
El m\'etodo de escorrent\'ia y el tiempo de concentraci\'on se entienden & S\'i / Parcial / No & \\
La delimitaci\'on de subcuencas y el evento de dise\~no parecen coherentes & S\'i / Parcial / No & \\
Los puntos cr\'iticos del sistema pluvial son identificables & S\'i / Parcial / No & \\
\bottomrule
\end{tabular}

\subsection*{5. Integraci\'on y presentaci\'on}
\begin{tabular}{p{0.54\linewidth} p{0.12\linewidth} p{0.20\linewidth}}
\toprule
\textbf{Criterio} & \textbf{Estado} & \textbf{Observaci\'on} \\
\midrule
Hay coherencia entre memoria, planos y resultados de modelo & S\'i / Parcial / No & \\
Las decisiones t\'ecnicas parecen justificadas & S\'i / Parcial / No & \\
El borrador permite llegar a una entrega final ordenada & S\'i / Parcial / No & \\
\bottomrule
\end{tabular}

\section{Hallazgos obligatorios}
\subsection*{Fortalezas del trabajo revisado}
\begin{itemize}
\item
\item
\end{itemize}

\subsection*{Observaciones prioritarias antes de la entrega final}
\begin{enumerate}
\item
\item
\item
\end{enumerate}

\subsection*{Sugerencias puntuales de mejora}
\begin{itemize}
\item
\item
\item
\end{itemize}

\section{Conclusi\'on de la revisi\'on}
\begin{tabular}{p{0.32\linewidth} p{0.58\linewidth}}
\toprule
Valoraci\'on general del borrador & S\'olido / Aceptable con ajustes / Requiere revisi\'on importante \\
Aspecto m\'as urgente a corregir & \\
Aspecto mejor resuelto del equipo revisado & \\
\bottomrule
\end{tabular}

\section{Constancia}
\begin{warnbox}
Este formulario debe entregarse como evidencia de revisi\'on por pares y, al mismo tiempo, compartirse con el equipo revisado para que pueda incorporar observaciones antes de la entrega final.
\end{warnbox}

\begin{tabular}{p{0.42\linewidth} p{0.42\linewidth}}
\toprule
Firma o nombre del equipo revisor & Recibido por el equipo revisado \\
\midrule
 & \\
 & \\
\bottomrule
\end{tabular}

\end{document}
